\chapter{Unit 3: Semiconductor PN Junction Diode Forward Characteristics \& Applications}

\section{Section A: Multiple Choice Practice Questions (60 Curated Problems)}

\noindent The following 60 multiple-choice questions reflect the official workshop evaluation bank, categorized by fundamental learning outcome and laboratory testing difficulty.\vspace{3mm}

\noindent\textbf{Q1. What does forward biasing a PN junction diode mean?} \hfill [Tag: Analysis of the voltage-current characteristics of a PN junction diode under forward-bias conditions | Easy]
\begin{enumerate}[label=(\alph*),topsep=2pt,itemsep=1pt,partopsep=0pt,parsep=0pt]
    \item {Connecting both sides to the positive terminal}
    \item {Connecting P-side to negative and N-side to positive}
    \item {Connecting both sides to the negative terminal}
    \item {Connecting P-side to positive and N-side to negative}
\end{enumerate}
\vspace{2mm}

\noindent\textbf{Q2. What happens to the depletion-region width when a PN junction diode is forward biased?} \hfill [Tag: Analysis of the voltage-current characteristics of a PN junction diode under forward-bias conditions | Easy]
\begin{enumerate}[label=(\alph*),topsep=2pt,itemsep=1pt,partopsep=0pt,parsep=0pt]
    \item {It becomes infinite}
    \item {It becomes narrower}
    \item {It becomes wider}
    \item {It remains unchanged}
\end{enumerate}
\vspace{2mm}

\noindent\textbf{Q3. What happens to the barrier potential of a PN junction under forward bias?} \hfill [Tag: Analysis of the voltage-current characteristics of a PN junction diode under forward-bias conditions | Easy]
\begin{enumerate}[label=(\alph*),topsep=2pt,itemsep=1pt,partopsep=0pt,parsep=0pt]
    \item {It remains constant}
    \item {It increases}
    \item {It decreases}
    \item {It reverses polarity}
\end{enumerate}
\vspace{2mm}

\noindent\textbf{Q4. Which charge carriers mainly contribute to current in a forward-biased PN junction diode?} \hfill [Tag: Analysis of the voltage-current characteristics of a PN junction diode under forward-bias conditions | Easy]
\begin{enumerate}[label=(\alph*),topsep=2pt,itemsep=1pt,partopsep=0pt,parsep=0pt]
    \item {Stationary ions}
    \item {Majority carriers}
    \item {Minority carriers}
    \item {Neutral atoms}
\end{enumerate}
\vspace{2mm}

\noindent\textbf{Q5. What is the typical cut-in voltage of a silicon diode?} \hfill [Tag: Analysis of the voltage-current characteristics of a PN junction diode under forward-bias conditions | Easy]
\begin{enumerate}[label=(\alph*),topsep=2pt,itemsep=1pt,partopsep=0pt,parsep=0pt]
    \item {2.0 V}
    \item {0.3 V}
    \item {0.7 V}
    \item {1.2 V}
\end{enumerate}
\vspace{2mm}

\noindent\textbf{Q6. What is the typical cut-in voltage of a germanium diode?} \hfill [Tag: Analysis of the voltage-current characteristics of a PN junction diode under forward-bias conditions | Easy]
\begin{enumerate}[label=(\alph*),topsep=2pt,itemsep=1pt,partopsep=0pt,parsep=0pt]
    \item {1.5 V}
    \item {1.0 V}
    \item {0.7 V}
    \item {0.3 V}
\end{enumerate}
\vspace{2mm}

\noindent\textbf{Q7. What is another common name for the cut-in voltage of a diode?} \hfill [Tag: Analysis of the voltage-current characteristics of a PN junction diode under forward-bias conditions | Easy]
\begin{enumerate}[label=(\alph*),topsep=2pt,itemsep=1pt,partopsep=0pt,parsep=0pt]
    \item {Reverse voltage}
    \item {Breakdown voltage}
    \item {Knee voltage}
    \item {Peak voltage}
\end{enumerate}
\vspace{2mm}

\noindent\textbf{Q8. How does forward current change after the diode voltage exceeds the knee voltage?} \hfill [Tag: Analysis of the voltage-current characteristics of a PN junction diode under forward-bias conditions | Easy]
\begin{enumerate}[label=(\alph*),topsep=2pt,itemsep=1pt,partopsep=0pt,parsep=0pt]
    \item {It rises rapidly}
    \item {It becomes zero}
    \item {It falls rapidly}
    \item {It stays constant}
\end{enumerate}
\vspace{2mm}

\noindent\textbf{Q9. Which axis normally represents diode current on a voltage-current characteristic graph?} \hfill [Tag: Analysis of the voltage-current characteristics of a PN junction diode under forward-bias conditions | Easy]
\begin{enumerate}[label=(\alph*),topsep=2pt,itemsep=1pt,partopsep=0pt,parsep=0pt]
    \item {Circular axis}
    \item {Diagonal axis}
    \item {Horizontal axis}
    \item {Vertical axis}
\end{enumerate}
\vspace{2mm}

\noindent\textbf{Q10. Which axis normally represents diode voltage on a voltage-current characteristic graph?} \hfill [Tag: Analysis of the voltage-current characteristics of a PN junction diode under forward-bias conditions | Easy]
\begin{enumerate}[label=(\alph*),topsep=2pt,itemsep=1pt,partopsep=0pt,parsep=0pt]
    \item {Horizontal axis}
    \item {Diagonal axis}
    \item {Vertical axis}
    \item {Circular axis}
\end{enumerate}
\vspace{2mm}

\noindent\textbf{Q11. In which quadrant is the forward-bias characteristic of a diode usually plotted?} \hfill [Tag: Analysis of the voltage-current characteristics of a PN junction diode under forward-bias conditions | Easy]
\begin{enumerate}[label=(\alph*),topsep=2pt,itemsep=1pt,partopsep=0pt,parsep=0pt]
    \item {Third quadrant}
    \item {Fourth quadrant}
    \item {First quadrant}
    \item {Second quadrant}
\end{enumerate}
\vspace{2mm}

\noindent\textbf{Q12. How does an ideal diode behave when it is forward biased?} \hfill [Tag: Analysis of the voltage-current characteristics of a PN junction diode under forward-bias conditions | Easy]
\begin{enumerate}[label=(\alph*),topsep=2pt,itemsep=1pt,partopsep=0pt,parsep=0pt]
    \item {As a closed switch}
    \item {As a current source}
    \item {As an open switch}
    \item {As a charged capacitor}
\end{enumerate}
\vspace{2mm}

\noindent\textbf{Q13. What is the forward-voltage drop across an ideal conducting diode?} \hfill [Tag: Analysis of the voltage-current characteristics of a PN junction diode under forward-bias conditions | Easy]
\begin{enumerate}[label=(\alph*),topsep=2pt,itemsep=1pt,partopsep=0pt,parsep=0pt]
    \item {0.3 V}
    \item {0 V}
    \item {0.7 V}
    \item {1.0 V}
\end{enumerate}
\vspace{2mm}

\noindent\textbf{Q14. Before the forward voltage reaches the knee voltage, the diode current is generally:} \hfill [Tag: Analysis of the voltage-current characteristics of a PN junction diode under forward-bias conditions | Easy]
\begin{enumerate}[label=(\alph*),topsep=2pt,itemsep=1pt,partopsep=0pt,parsep=0pt]
    \item {Very large}
    \item {Very small}
    \item {Exactly infinite}
    \item {Constantly negative}
\end{enumerate}
\vspace{2mm}

\noindent\textbf{Q15. Which instrument is connected in series to measure the forward current of a diode?} \hfill [Tag: Analysis of the voltage-current characteristics of a PN junction diode under forward-bias conditions | Easy]
\begin{enumerate}[label=(\alph*),topsep=2pt,itemsep=1pt,partopsep=0pt,parsep=0pt]
    \item {Wattmeter}
    \item {Ammeter}
    \item {Voltmeter}
    \item {Ohmmeter}
\end{enumerate}
\vspace{2mm}

\noindent\textbf{Q16. Which instrument is connected across a diode to measure its forward voltage?} \hfill [Tag: Analysis of the voltage-current characteristics of a PN junction diode under forward-bias conditions | Easy]
\begin{enumerate}[label=(\alph*),topsep=2pt,itemsep=1pt,partopsep=0pt,parsep=0pt]
    \item {Ammeter}
    \item {Voltmeter}
    \item {Galvanometer}
    \item {Frequency meter}
\end{enumerate}
\vspace{2mm}

\noindent\textbf{Q17. Why is a resistor commonly connected in series with a forward-biased diode?} \hfill [Tag: Analysis of the voltage-current characteristics of a PN junction diode under forward-bias conditions | Easy]
\begin{enumerate}[label=(\alph*),topsep=2pt,itemsep=1pt,partopsep=0pt,parsep=0pt]
    \item {To create breakdown}
    \item {To limit current}
    \item {To increase voltage}
    \item {To reverse polarity}
\end{enumerate}
\vspace{2mm}

\noindent\textbf{Q18. The static forward resistance of a diode at an operating point is calculated using which expression?} \hfill [Tag: Analysis of the voltage-current characteristics of a PN junction diode under forward-bias conditions | Easy]
\begin{enumerate}[label=(\alph*),topsep=2pt,itemsep=1pt,partopsep=0pt,parsep=0pt]
    \item {$R_F = V_F \cdot I_F$}
    \item {$R_F = V_F + I_F$}
    \item {$R_F = \frac{I_F}{V_F}$}
    \item {$R_F = \frac{V_F}{I_F}$}
\end{enumerate}
\vspace{2mm}

\noindent\textbf{Q19. A diode has a forward voltage of 0.6 V and a forward current of 3 mA . What is its static forward resistance?} \hfill [Tag: Analysis of the voltage-current characteristics of a PN junction diode under forward-bias conditions | Easy]
\begin{enumerate}[label=(\alph*),topsep=2pt,itemsep=1pt,partopsep=0pt,parsep=0pt]
    \item {200 $\Omega$}
    \item {300 $\Omega$}
    \item {600 $\Omega$}
    \item {100 $\Omega$}
\end{enumerate}
\vspace{2mm}

\noindent\textbf{Q20. What does the nonlinear shape of a diode's forward characteristic indicate?} \hfill [Tag: Analysis of the voltage-current characteristics of a PN junction diode under forward-bias conditions | Easy]
\begin{enumerate}[label=(\alph*),topsep=2pt,itemsep=1pt,partopsep=0pt,parsep=0pt]
    \item {Resistance remains perfectly constant}
    \item {Current is not directly proportional to voltage}
    \item {Current is always equal to voltage}
    \item {Voltage remains zero at all currents}
\end{enumerate}
\vspace{2mm}

\noindent\textbf{Q21. A silicon PN junction diode is forward biased. At which voltage does its current typically begin to increase significantly?} \hfill [Tag: Analysis of the voltage-current characteristics of a PN junction diode under forward-bias conditions | Medium]
\begin{enumerate}[label=(\alph*),topsep=2pt,itemsep=1pt,partopsep=0pt,parsep=0pt]
    \item {About 0.5 V}
    \item {About 0.7 V}
    \item {About 1.2 V}
    \item {About 0.3 V}
\end{enumerate}
\vspace{2mm}

\noindent\textbf{Q22. For a silicon diode operating near its knee voltage, what is the most likely effect of increasing the forward voltage from 0.7 V to 0.8 V ?} \hfill [Tag: Analysis of the voltage-current characteristics of a PN junction diode under forward-bias conditions | Medium]
\begin{enumerate}[label=(\alph*),topsep=2pt,itemsep=1pt,partopsep=0pt,parsep=0pt]
    \item {The diode becomes reverse biased}
    \item {The current decreases slightly}
    \item {The current increases substantially}
    \item {The current remains exactly constant}
\end{enumerate}
\vspace{2mm}

\noindent\textbf{Q23. The forward current of a diode changes from 2 mA to 3 mA when its forward voltage changes from 0.68 V to 0.70 V . What is its approximate dc or small-signal resistance over this interval?} \hfill [Tag: Analysis of the voltage-current characteristics of a PN junction diode under forward-bias conditions | Medium]
\begin{enumerate}[label=(\alph*),topsep=2pt,itemsep=1pt,partopsep=0pt,parsep=0pt]
    \item {200 $\Omega$}
    \item {50 $\Omega$}
    \item {5 $\Omega$}
    \item {20 $\Omega$}
\end{enumerate}
\vspace{2mm}

\noindent\textbf{Q24. A diode has a forward voltage of 0.72 V and carries a current of 15 mA . What power is dissipated in the diode?} \hfill [Tag: Analysis of the voltage-current characteristics of a PN junction diode under forward-bias conditions | Medium]
\begin{enumerate}[label=(\alph*),topsep=2pt,itemsep=1pt,partopsep=0pt,parsep=0pt]
    \item {10.8 mW}
    \item {48 mW}
    \item {1.08 mW}
    \item {108 mW}
\end{enumerate}
\vspace{2mm}

\noindent\textbf{Q25. When the temperature of a silicon diode increases, how does its forward characteristic generally change at a fixed forward current?} \hfill [Tag: Analysis of the voltage-current characteristics of a PN junction diode under forward-bias conditions | Medium]
\begin{enumerate}[label=(\alph*),topsep=2pt,itemsep=1pt,partopsep=0pt,parsep=0pt]
    \item {The characteristic becomes vertical}
    \item {The required voltage increases}
    \item {The required voltage decreases}
    \item {The current becomes zero}
\end{enumerate}
\vspace{2mm}

\noindent\textbf{Q26. A 5 V source is connected in series with a 1 k $\Omega$ resistor and a silicon diode modeled with a constant forward drop of 0.7 V . What is the approximate diode current?} \hfill [Tag: Analysis of the voltage-current characteristics of a PN junction diode under forward-bias conditions | Medium]
\begin{enumerate}[label=(\alph*),topsep=2pt,itemsep=1pt,partopsep=0pt,parsep=0pt]
    \item {0.43 mA}
    \item {4.3 mA}
    \item {7.0 mA}
    \item {5.7 mA}
\end{enumerate}
\vspace{2mm}

\noindent\textbf{Q27. On a forward-bias characteristic curve, the current changes from 4 mA to 6 mA while the voltage changes from 0.64 V to 0.66 V . What is the approximate slope $\Delta$ I /$\Delta$ V ?} \hfill [Tag: Analysis of the voltage-current characteristics of a PN junction diode under forward-bias conditions | Medium]
\begin{enumerate}[label=(\alph*),topsep=2pt,itemsep=1pt,partopsep=0pt,parsep=0pt]
    \item {10 S}
    \item {100 S}
    \item {0.1 S}
    \item {50 S}
\end{enumerate}
\vspace{2mm}

\noindent\textbf{Q28. Which feature best identifies the knee region of a PN junction diode's forward characteristic?} \hfill [Tag: Analysis of the voltage-current characteristics of a PN junction diode under forward-bias conditions | Medium]
\begin{enumerate}[label=(\alph*),topsep=2pt,itemsep=1pt,partopsep=0pt,parsep=0pt]
    \item {Current changes rapidly after a small voltage rise}
    \item {Voltage and current both remain constant}
    \item {Voltage changes rapidly after a small current rise}
    \item {Current remains zero for every voltage}
\end{enumerate}
\vspace{2mm}

\noindent\textbf{Q29. Compared with a silicon diode, a germanium diode generally shows which forward-bias characteristic?} \hfill [Tag: Analysis of the voltage-current characteristics of a PN junction diode under forward-bias conditions | Medium]
\begin{enumerate}[label=(\alph*),topsep=2pt,itemsep=1pt,partopsep=0pt,parsep=0pt]
    \item {A higher cut-in voltage}
    \item {An identical cut-in voltage}
    \item {A lower cut-in voltage}
    \item {No forward current}
\end{enumerate}
\vspace{2mm}

\noindent\textbf{Q30. Using the diode equation with the same saturation current and ideality factor, the forward voltage increases from 0.60 V to 0.65 V at room temperature. Approximately how does the current change for n = 1 ?} \hfill [Tag: Analysis of the voltage-current characteristics of a PN junction diode under forward-bias conditions | Medium]
\begin{enumerate}[label=(\alph*),topsep=2pt,itemsep=1pt,partopsep=0pt,parsep=0pt]
    \item {It becomes about half}
    \item {It remains nearly unchanged}
    \item {It increases by exactly 25 times}
    \item {It increases by about 7.4 times}
\end{enumerate}
\vspace{2mm}

\noindent\textbf{Q31. For a forward-biased diode, which region of the characteristic is most accurately modeled by an approximately constant voltage drop?} \hfill [Tag: Analysis of the voltage-current characteristics of a PN junction diode under forward-bias conditions | Medium]
\begin{enumerate}[label=(\alph*),topsep=2pt,itemsep=1pt,partopsep=0pt,parsep=0pt]
    \item {The region after the knee voltage}
    \item {The region near reverse breakdown}
    \item {The region before any applied voltage}
    \item {The region below zero current}
\end{enumerate}
\vspace{2mm}

\noindent\textbf{Q32. If the forward voltage of a diode is plotted on the horizontal axis and forward current on the vertical axis, what does a steep curve indicate?} \hfill [Tag: Analysis of the voltage-current characteristics of a PN junction diode under forward-bias conditions | Medium]
\begin{enumerate}[label=(\alph*),topsep=2pt,itemsep=1pt,partopsep=0pt,parsep=0pt]
    \item {A small dynamic resistance}
    \item {A zero saturation current}
    \item {A reverse-biased condition}
    \item {A large dynamic resistance}
\end{enumerate}
\vspace{2mm}

\noindent\textbf{Q33. For a diode with ideality factor n = 2 at room temperature, approximately what forward-voltage increase is required to increase the current by a factor of ten?} \hfill [Tag: Analysis of the voltage-current characteristics of a PN junction diode under forward-bias conditions | Medium]
\begin{enumerate}[label=(\alph*),topsep=2pt,itemsep=1pt,partopsep=0pt,parsep=0pt]
    \item {120 mV}
    \item {240 mV}
    \item {60 mV}
    \item {30 mV}
\end{enumerate}
\vspace{2mm}

\noindent\textbf{Q34. A measured diode characteristic shows a nearly straight line with a nonzero slope after the knee region. What does the slope mainly represent?} \hfill [Tag: Analysis of the voltage-current characteristics of a PN junction diode under forward-bias conditions | Medium]
\begin{enumerate}[label=(\alph*),topsep=2pt,itemsep=1pt,partopsep=0pt,parsep=0pt]
    \item {The diode's reverse voltage}
    \item {The diode's series resistance}
    \item {The diode's junction capacitance}
    \item {The diode's breakdown current}
\end{enumerate}
\vspace{2mm}

\noindent\textbf{Q35. At room temperature, a diode carries 10 mA and has ideality factor n = 1 . Using r d $\approx$ n V T / I D , what is its approximate dynamic resistance?} \hfill [Tag: Analysis of the voltage-current characteristics of a PN junction diode under forward-bias conditions | Medium]
\begin{enumerate}[label=(\alph*),topsep=2pt,itemsep=1pt,partopsep=0pt,parsep=0pt]
    \item {25 $\Omega$}
    \item {2.5 $\Omega$}
    \item {0.25 $\Omega$}
    \item {250 $\Omega$}
\end{enumerate}
\vspace{2mm}

\noindent\textbf{Q36. A forward-biased diode has a current of 1 mA at 0.60 V and 10 mA at 0.66 V . Which conclusion is most reasonable at room temperature?} \hfill [Tag: Analysis of the voltage-current characteristics of a PN junction diode under forward-bias conditions | Medium]
\begin{enumerate}[label=(\alph*),topsep=2pt,itemsep=1pt,partopsep=0pt,parsep=0pt]
    \item {The ideality factor is close to 3}
    \item {The ideality factor is close to 1}
    \item {The diode has zero saturation current}
    \item {The diode is reverse biased}
\end{enumerate}
\vspace{2mm}

\noindent\textbf{Q37. A diode and a resistor are connected to a dc source. If the source voltage increases while the diode remains in forward conduction, what primarily limits the increase in diode current?} \hfill [Tag: Analysis of the voltage-current characteristics of a PN junction diode under forward-bias conditions | Medium]
\begin{enumerate}[label=(\alph*),topsep=2pt,itemsep=1pt,partopsep=0pt,parsep=0pt]
    \item {The depletion layer alone}
    \item {The reverse saturation current}
    \item {The absence of forward bias}
    \item {The resistor in series}
\end{enumerate}
\vspace{2mm}

\noindent\textbf{Q38. For a forward-biased silicon diode, the current is 5 mA at 2 5 $^\circ$ C . If the temperature rises while the current is held constant, what happens to the forward voltage?} \hfill [Tag: Analysis of the voltage-current characteristics of a PN junction diode under forward-bias conditions | Medium]
\begin{enumerate}[label=(\alph*),topsep=2pt,itemsep=1pt,partopsep=0pt,parsep=0pt]
    \item {It reverses polarity immediately}
    \item {It remains exactly unchanged}
    \item {It increases by about 2 mV / $^\circ$ C}
    \item {It decreases by about 2 mV / $^\circ$ C}
\end{enumerate}
\vspace{2mm}

\noindent\textbf{Q39. A diode's forward characteristic is initially almost horizontal and then becomes steep. What physical interpretation best matches the steep portion?} \hfill [Tag: Analysis of the voltage-current characteristics of a PN junction diode under forward-bias conditions | Medium]
\begin{enumerate}[label=(\alph*),topsep=2pt,itemsep=1pt,partopsep=0pt,parsep=0pt]
    \item {The diode current is independent of voltage}
    \item {A small voltage change produces a large current change}
    \item {The diode is operating in reverse breakdown}
    \item {The junction barrier is increasing strongly}
\end{enumerate}
\vspace{2mm}

\noindent\textbf{Q40. A load line intersects a diode's forward I - V characteristic at 0.68 V and 8 mA . What does this intersection represent?} \hfill [Tag: Analysis of the voltage-current characteristics of a PN junction diode under forward-bias conditions | Medium]
\begin{enumerate}[label=(\alph*),topsep=2pt,itemsep=1pt,partopsep=0pt,parsep=0pt]
    \item {The diode's reverse breakdown point}
    \item {The diode's zero-current voltage}
    \item {The circuit's forward operating point}
    \item {The maximum possible source voltage}
\end{enumerate}
\vspace{2mm}

\noindent\textbf{Q41. At 300 K , a diode carries 0.10 mA at 0.50 V and 1.00 mA at 0.62 V . Assuming the Shockley model and negligible series resistance, what is the ideality factor n ? Use V T = 25.85 mV .} \hfill [Tag: Analysis of the voltage-current characteristics of a PN junction diode under forward-bias conditions | Hard]
\begin{enumerate}[label=(\alph*),topsep=2pt,itemsep=1pt,partopsep=0pt,parsep=0pt]
    \item {n $\approx$ 2.02}
    \item {n $\approx$ 1.01}
    \item {n $\approx$ 1.52}
    \item {n $\approx$ 2.51}
\end{enumerate}
\vspace{2mm}

\noindent\textbf{Q42. A diode at 300 K carries 1.00 mA at 0.600 V . If n = 1 and the - 1 term is negligible, what is its reverse saturation current?} \hfill [Tag: Analysis of the voltage-current characteristics of a PN junction diode under forward-bias conditions | Hard]
\begin{enumerate}[label=(\alph*),topsep=2pt,itemsep=1pt,partopsep=0pt,parsep=0pt]
    \item {2.41 $\times$ 1 0 - 13 A}
    \item {1.20 $\times$ 1 0 - 12 A}
    \item {8.30 $\times$ 1 0 - 14 A}
    \item {4.15 $\times$ 1 0 - 14 A}
\end{enumerate}
\vspace{2mm}

\noindent\textbf{Q43. A forward-biased diode operates at I D = 2.00 mA with n = 2 at 300 K . Neglecting series resistance and assuming I D $\gg$ I S , what is its small-signal resistance?} \hfill [Tag: Analysis of the voltage-current characteristics of a PN junction diode under forward-bias conditions | Hard]
\begin{enumerate}[label=(\alph*),topsep=2pt,itemsep=1pt,partopsep=0pt,parsep=0pt]
    \item {51.70 $\Omega$}
    \item {25.85 $\Omega$}
    \item {103.4 $\Omega$}
    \item {12.93 $\Omega$}
\end{enumerate}
\vspace{2mm}

\noindent\textbf{Q44. A diode with n = 1 and V T = 25.0 mV is biased at 1.00 mA . Its voltage is changed by + 25.0 mV and then by - 25.0 mV relative to the bias point. What is the ratio of the upward current change to the magnitude of the downward current change?} \hfill [Tag: Analysis of the voltage-current characteristics of a PN junction diode under forward-bias conditions | Hard]
\begin{enumerate}[label=(\alph*),topsep=2pt,itemsep=1pt,partopsep=0pt,parsep=0pt]
    \item {2.718}
    \item {7.389}
    \item {1.718}
    \item {1.000}
\end{enumerate}
\vspace{2mm}

\noindent\textbf{Q45. A diode with I S = 1 0 - 12 A and n = 1 is connected in series with 1.00 k $\Omega$ across a 5.00 V supply. At 300 K , which operating point most closely satisfies both the diode equation and the load line?} \hfill [Tag: Analysis of the voltage-current characteristics of a PN junction diode under forward-bias conditions | Hard]
\begin{enumerate}[label=(\alph*),topsep=2pt,itemsep=1pt,partopsep=0pt,parsep=0pt]
    \item {I D $\approx$ 4.70 mA , V D $\approx$ 0.300 V}
    \item {I D $\approx$ 4.00 mA , V D $\approx$ 1.00 V}
    \item {I D $\approx$ 4.43 mA , V D $\approx$ 0.574 V}
    \item {I D $\approx$ 5.00 mA , V D $\approx$ 0.535 V}
\end{enumerate}
\vspace{2mm}

\noindent\textbf{Q46. At 300 K , a diode with n = 1 has measured terminal voltages of 0.750 V at 10.0 mA and 1.000 V at 100 mA . Assuming a constant series resistance, estimate that resistance.} \hfill [Tag: Analysis of the voltage-current characteristics of a PN junction diode under forward-bias conditions | Hard]
\begin{enumerate}[label=(\alph*),topsep=2pt,itemsep=1pt,partopsep=0pt,parsep=0pt]
    \item {2.78 $\Omega$}
    \item {2.12 $\Omega$}
    \item {1.06 $\Omega$}
    \item {4.24 $\Omega$}
\end{enumerate}
\vspace{2mm}

\noindent\textbf{Q47. A diode carries a fixed current at 300 K with a forward voltage of 0.650 V . Its saturation current doubles for every 8 K rise, while n = 1 remains constant. According to the Shockley model, what is the forward voltage at 340 K ? Use V T = 25.85 mV at 300 K .} \hfill [Tag: Analysis of the voltage-current characteristics of a PN junction diode under forward-bias conditions | Hard]
\begin{enumerate}[label=(\alph*),topsep=2pt,itemsep=1pt,partopsep=0pt,parsep=0pt]
    \item {0.560 V}
    \item {0.737 V}
    \item {0.600 V}
    \item {0.635 V}
\end{enumerate}
\vspace{2mm}

\noindent\textbf{Q48. The straight-line region of a diode's forward characteristic plotted as V D versus lo g 10 I D has a slope of 89.3 mV per decade at 300 K . What ideality factor does this imply?} \hfill [Tag: Analysis of the voltage-current characteristics of a PN junction diode under forward-bias conditions | Hard]
\begin{enumerate}[label=(\alph*),topsep=2pt,itemsep=1pt,partopsep=0pt,parsep=0pt]
    \item {n $\approx$ 1.00}
    \item {n $\approx$ 1.25}
    \item {n $\approx$ 1.50}
    \item {n $\approx$ 2.00}
\end{enumerate}
\vspace{2mm}

\noindent\textbf{Q49. A practical diode has n = 2 , series resistance R s = 3.00 $\Omega$ , and a bias current of 10.0 mA at 300 K . What is its terminal small-signal conductance?} \hfill [Tag: Analysis of the voltage-current characteristics of a PN junction diode under forward-bias conditions | Hard]
\begin{enumerate}[label=(\alph*),topsep=2pt,itemsep=1pt,partopsep=0pt,parsep=0pt]
    \item {0.0785 S}
    \item {0.122 S}
    \item {0.193 S}
    \item {0.333 S}
\end{enumerate}
\vspace{2mm}

\noindent\textbf{Q50. A diode with n = 1 is biased at 2.00 mA through a 100 $\Omega$ resistor. At 300 K , what is the small-signal change in diode current per volt of supply variation, assuming no diode series resistance?} \hfill [Tag: Analysis of the voltage-current characteristics of a PN junction diode under forward-bias conditions | Hard]
\begin{enumerate}[label=(\alph*),topsep=2pt,itemsep=1pt,partopsep=0pt,parsep=0pt]
    \item {8.86 mA/V}
    \item {77.4 mA/V}
    \item {10.0 mA/V}
    \item {7.97 mA/V}
\end{enumerate}
\vspace{2mm}

\noindent\textbf{Q51. For a diode with I S = 10.0 nA , n = 2 , and V T = 25.85 mV , what current flows at a small forward voltage of 10.0 mV when the full Shockley equation is used?} \hfill [Tag: Analysis of the voltage-current characteristics of a PN junction diode under forward-bias conditions | Hard]
\begin{enumerate}[label=(\alph*),topsep=2pt,itemsep=1pt,partopsep=0pt,parsep=0pt]
    \item {2.13 nA}
    \item {10.0 nA}
    \item {8.24 nA}
    \item {12.1 nA}
\end{enumerate}
\vspace{2mm}

\noindent\textbf{Q52. A diode operates at V D = 0.700 V and I D = 10.0 mA with n = 2 at 300 K . Neglecting series resistance, which pair gives its DC resistance and small-signal resistance, respectively?} \hfill [Tag: Analysis of the voltage-current characteristics of a PN junction diode under forward-bias conditions | Hard]
\begin{enumerate}[label=(\alph*),topsep=2pt,itemsep=1pt,partopsep=0pt,parsep=0pt]
    \item {35.0 $\Omega$ and 5.17 $\Omega$}
    \item {5.17 $\Omega$ and 70.0 $\Omega$}
    \item {70.0 $\Omega$ and 5.17 $\Omega$}
    \item {70.0 $\Omega$ and 2.59 $\Omega$}
\end{enumerate}
\vspace{2mm}

\noindent\textbf{Q53. At 300 K , the forward current of a PN diode doubles whenever its voltage increases by approximately 36.0 mV over a certain low-current range. Which interpretation is most consistent with this observation?} \hfill [Tag: Analysis of the voltage-current characteristics of a PN junction diode under forward-bias conditions | Hard]
\begin{enumerate}[label=(\alph*),topsep=2pt,itemsep=1pt,partopsep=0pt,parsep=0pt]
    \item {n $\approx$ 2 , suggesting recombination-dominated current}
    \item {n $\approx$ 1 , suggesting diffusion-dominated current}
    \item {n $\approx$ 4 , suggesting series-resistance domination}
    \item {n $\approx$ 0.5 , suggesting tunneling-dominated current}
\end{enumerate}
\vspace{2mm}

\noindent\textbf{Q54. For a diode described by V D = n V T ln ( I D / I S ) + I D R s , the measured values of d V D / d ln I D are 45.85 mV at 10.0 mA and 125.85 mV at 50.0 mA . At 300 K , determine n and R s .} \hfill [Tag: Analysis of the voltage-current characteristics of a PN junction diode under forward-bias conditions | Hard]
\begin{enumerate}[label=(\alph*),topsep=2pt,itemsep=1pt,partopsep=0pt,parsep=0pt]
    \item {n = 1.0 and R s = 2.0 $\Omega$}
    \item {n = 2.0 and R s = 2.0 $\Omega$}
    \item {n = 2.0 and R s = 1.0 $\Omega$}
    \item {n = 1.0 and R s = 4.0 $\Omega$}
\end{enumerate}
\vspace{2mm}

\noindent\textbf{Q55. Two parallel diodes have equal ideality factors and temperatures but saturation currents satisfying I S 1 = 4 I S 2 . If their total forward current is 10.0 mA and series resistances are negligible, how is the current shared?} \hfill [Tag: Analysis of the voltage-current characteristics of a PN junction diode under forward-bias conditions | Hard]
\begin{enumerate}[label=(\alph*),topsep=2pt,itemsep=1pt,partopsep=0pt,parsep=0pt]
    \item {I 1 = 8.0 mA , I 2 = 2.0 mA}
    \item {I 1 = 5.0 mA , I 2 = 5.0 mA}
    \item {I 1 = 9.0 mA , I 2 = 1.0 mA}
    \item {I 1 = 6.0 mA , I 2 = 4.0 mA}
\end{enumerate}
\vspace{2mm}

\noindent\textbf{Q56. Two parallel diodes each initially carry 1.00 mA at the same voltage. Their ideality factors are n 1 = 1 and n 2 = 2 . At 300 K , the common voltage is increased by 51.7 mV . Neglecting series resistance, what is the new total current?} \hfill [Tag: Analysis of the voltage-current characteristics of a PN junction diode under forward-bias conditions | Hard]
\begin{enumerate}[label=(\alph*),topsep=2pt,itemsep=1pt,partopsep=0pt,parsep=0pt]
    \item {7.39 mA}
    \item {14.78 mA}
    \item {5.44 mA}
    \item {10.11 mA}
\end{enumerate}
\vspace{2mm}

\noindent\textbf{Q57. A diode's terminal voltage rises by 160 mV as its current increases from 1.00 mA to 10.0 mA . If its series resistance is 10.0 $\Omega$ , what ideality factor is implied at 300 K ?} \hfill [Tag: Analysis of the voltage-current characteristics of a PN junction diode under forward-bias conditions | Hard]
\begin{enumerate}[label=(\alph*),topsep=2pt,itemsep=1pt,partopsep=0pt,parsep=0pt]
    \item {n $\approx$ 1.68}
    \item {n $\approx$ 2.69}
    \item {n $\approx$ 0.84}
    \item {n $\approx$ 1.18}
\end{enumerate}
\vspace{2mm}

\noindent\textbf{Q58. A diode has n = 1.5 and an uncertain saturation current as high as 5.0 $\times$ 1 0 - 12 A . At 300 K , what is the maximum applied forward voltage that guarantees the Shockley current does not exceed 20.0 mA ? Neglect series resistance.} \hfill [Tag: Analysis of the voltage-current characteristics of a PN junction diode under forward-bias conditions | Hard]
\begin{enumerate}[label=(\alph*),topsep=2pt,itemsep=1pt,partopsep=0pt,parsep=0pt]
    \item {0.786 V}
    \item {0.824 V}
    \item {0.857 V}
    \item {0.914 V}
\end{enumerate}
\vspace{2mm}

\noindent\textbf{Q59. An ideal diode with n = 1 and V T = 25.0 mV is biased at 1.00 mA . A zero-average sinusoidal voltage v ( t ) = 10.0 sin ( $\omega$ t ) mV is superimposed. Neglecting I S relative to forward current, what is the resulting average current?} \hfill [Tag: Analysis of the voltage-current characteristics of a PN junction diode under forward-bias conditions | Hard]
\begin{enumerate}[label=(\alph*),topsep=2pt,itemsep=1pt,partopsep=0pt,parsep=0pt]
    \item {1.173 mA}
    \item {1.040 mA}
    \item {1.000 mA}
    \item {1.020 mA}
\end{enumerate}
\vspace{2mm}

\noindent\textbf{Q60. A diode carries 1.00 mA at both 300 K and 350 K . Its ideality factor is 1 , while I S changes from 1 0 - 12 A at 300 K to 1 0 - 9 A at 350 K . Using V T = 25.85 mV and 30.16 mV , respectively, what is the forward-voltage change?} \hfill [Tag: Analysis of the voltage-current characteristics of a PN junction diode under forward-bias conditions | Hard]
\begin{enumerate}[label=(\alph*),topsep=2pt,itemsep=1pt,partopsep=0pt,parsep=0pt]
    \item {It decreases by approximately 119 mV}
    \item {It decreases by approximately 60 mV}
    \item {It increases by approximately 60 mV}
    \item {It increases by approximately 119 mV}
\end{enumerate}
\vspace{2mm}


\newpage
\section{Section B: Comprehensive Subjective Problem with Detailed Analytical Rubric}

\begin{subjectiveq}{3}{1}
A Silicon PN junction diode at room temperature ($T = 300\,\mathrm{K}$, $V_T = 26\,\mathrm{mV}$) has a reverse saturation current $I_s = 2\,\mathrm{nA}$ and ideality factor $\eta = 1$. (a) Determine the diode forward current when $V_D = 0.65\,\mathrm{V}$. (b) Calculate the static (DC) resistance and dynamic (AC) resistance at this operating point. (c) Explain why dynamic resistance is vastly smaller than static resistance in the forward conduction region.
\end{subjectiveq}

\begin{solbox}
\textbf{(a) Forward Conduction Current:}\\
Using the Shockley diode equation:\\
\[ I_D = I_s \left( e^{\frac{V_D}{\eta V_T}} - 1 \right) \]
Since $V_D / \eta V_T = 0.65 / 0.026 = 25 \gg 1$, the $-1$ term is negligible:\\
\[ I_D \approx 2 \times 10^{-9} \times e^{25} = 2 \times 10^{-9} \times 7.20 \times 10^{10} = 144.0\,\mathrm{mA} \]
\textbf{(b) Static and Dynamic Resistances:}\\
1. Static (DC) Resistance:\\
\[ R_D = \frac{V_D}{I_D} = \frac{0.65\,\mathrm{V}}{0.144\,\mathrm{A}} \approx 4.51\,\Omega \]
2. Dynamic (AC) Resistance:\\
\[ r_d = \frac{d V_D}{d I_D} = \frac{\eta V_T}{I_D} = \frac{0.026\,\mathrm{V}}{0.144\,\mathrm{A}} \approx 0.18\,\Omega \]
\textbf{(c) Physical Interpretation:}\\
The static resistance $R_D$ is the reciprocal of the slope of the secant line from the origin $(0,0)$ to the Q-point $(V_D, I_D)$. The dynamic resistance $r_d$ is the reciprocal of the tangent slope at the Q-point. Due to the exponential nature of forward conduction, the slope $d I_D / d V_D = I_D / (\eta V_T)$ increases steeply with current, causing the incremental resistance $r_d$ to become much smaller than $R_D$.
\end{solbox}

\begin{examinertip}
Examiner Rubric (10 Marks):\\
- Application of Shockley diode equation with correct exponential evaluation: 4 Marks\\
- Calculation of static resistance $R_D$: 2 Marks\\
- Derivation and numerical calculation of dynamic resistance $r_d = \eta V_T / I_D$: 2 Marks\\
- Graphical distinction between tangent and secant slopes: 2 Marks
\end{examinertip}


\section{Section C: Laboratory Viva-Voce Questions with Authoritative Solutions}

\noindent\textbf{Viva Question:} What is the knee or cut-in voltage of a diode, and what causes it physically?\\
\textbf{Technical Answer:} The cut-in voltage ($V_k \approx 0.7\text{ V}$ for Si, $0.3\text{ V}$ for Ge) is the forward bias potential required to counteract the built-in electrostatic potential barrier of the depletion region, allowing majority carriers to diffuse rapidly across the junction.
\vspace{3mm}

\noindent\textbf{Viva Question:} Why is a series current-limiting resistor mandatory during forward-bias diode testing?\\
\textbf{Technical Answer:} In the forward bias conduction region beyond cut-in, the diode dynamic resistance drops below $1\,\Omega$. Connecting a DC power supply directly across the diode without a series resistor leads to excessive current exceeding maximum forward ratings ($I_{F,\max}$), causing immediate burnout.
\vspace{3mm}

\noindent\textbf{Viva Question:} How do you test a diode using a Digital Multimeter's Diode Check function?\\
\textbf{Technical Answer:} In forward bias (red lead to anode, black to cathode), the DMM injects $\approx 1\text{ mA}$ and displays the forward drop ($0.55\text{--}0.75\text{ V}$ for Silicon). In reverse bias (red lead to cathode, black to anode), the display shows 'OL' (overload/infinite resistance). If it reads zero in both directions, the diode is shorted; if 'OL' in both, it is open.
\vspace{3mm}


\section{Section D: Workshop Troubleshooting \& Fault Diagnostic Case Study}

\textbf{Troubleshooting Scenario: Diode Circuit Emits Smoke in Rectifier Lab}\\
During a bridge rectifier experiment, an engineer switches on the transformer secondary ($15\\text{V}_{\\mathrm{rms}}$). Immediately, one 1N4007 diode overheats and cracks.\\\\
\textbf{Diagnostic Cause:} Diode polarity reversal. One diode was inserted backwards in the bridge configuration. When that half of the AC cycle occurred, the reversed diode and its counterpart created a direct dead short-circuit across the transformer secondary winding with zero series load resistance, drawing tens of amperes.\\
\textbf{Preventative Workshop Rule:} Always double check the cathode band (silver or white stripe) against the circuit schematic prior to energizing any power supply.
\vspace{4mm}
